\chapter{Logic}
\label{chap:logic}

\section*{Final review orientation}

本章按 source 重新整理 Lecture Notes Chapter 1 Logic、Lecture Notes Exercises 1--10、Worksheet 1、Homework 1、Homework 2。目标是把三件事练熟：第一，把 English statements 翻译成 Boolean formula 或 predicate formula；第二，用 truth table 和 logical equivalence 判断 deduction 是否 valid；第三，在 predicate logic 中正确处理 domain of discourse、quantifier order、negation 和 counterexample。

\section{Propositions and Boolean formulas}

Lecture Notes 把 mathematical argument 写成
\[
\begin{array}{c}
\text{Hypothesis A}\\
\text{Hypothesis B}\\
\hline
\text{Conclusion C}
\end{array}
\]
意思是：如果 hypotheses 都 true，那么 conclusion 必须 true。A proposition 是可以判断 true 或 false 的 sentence 或 claim。比如 ``The moon is made of cheese'' 是 proposition，虽然它是 false；``Close the door'' 不是 proposition，因为它是 command。

在 propositional logic 中，每个 proposition 被看作只取 \(T,F\) 两个 truth values 的 Boolean variable。主要 connectives 是
\[
\land,\qquad \lor,\qquad \lnot,\qquad \to,\qquad \leftrightarrow .
\]
\(A\to B\) 只有在 \(A=T,B=F\) 时 false；其他三种 rows 都 true。Boolean formula 必须 well-formed，且最终答案应使用 parentheses 消除 ambiguity。Lecture Notes 的读取顺序是先 \(\lnot\)，再 \(\land,\lor\)，最后 \(\to,\leftrightarrow\)。

\begin{exerciseblock}[Lecture Notes Exercise 1]
Lecture Notes 的 murder example 使用 variables \(A,\ldots,H\)。写出若干 logical dependencies。
\end{exerciseblock}

\begin{solution}
一种 source-consistent 的写法如下：
\[
\lnot A\to B,
\]
表示如果 murder 不是 during the night 发生，那么它 during the day 发生。
\[
(B\land D)\to \lnot G,
\]
表示如果 murder happened during the day，并且 Michael was asleep during the day，那么 Michael is not the murderer。
\[
(A\land C)\to \lnot F,
\]
表示如果 murder happened during the night，并且 Omar was asleep during the night，那么 Omar is not the murderer。
\[
E\to \lnot H,
\]
表示如果 Mykola 在相关时间段 asleep，那么 Mykola is not the murderer。
若额外假设 suspects 只可能是 Omar, Michael, Mykola，且至少一个人 guilty，可写
\[
F\lor G\lor H.
\]
若额外假设 exactly one murderer，还应加入 pairwise exclusion，例如
\[
(F\to\lnot G)\land(F\to\lnot H)\land(G\to\lnot F)\land(G\to\lnot H)\land(H\to\lnot F)\land(H\to\lnot G).
\]
\end{solution}

\begin{explanation}
Exercise 1 是 open-ended。关键不是猜唯一答案，而是每个 formula 都要对应一个明确 English assumption。最后两个 formulas 使用额外 background assumptions，写答案时必须说明。
\end{explanation}

\section{Truth tables and logical equivalence}

Basic truth tables 是本章所有判断的底层定义：
\[
\begin{array}{c|c}
A&\lnot A\\
\hline
F&T\\
T&F
\end{array}
\qquad
\begin{array}{cc|cc}
A&B&A\land B&A\lor B\\
\hline
F&F&F&F\\
F&T&F&T\\
T&F&F&T\\
T&T&T&T
\end{array}
\]
\[
\begin{array}{cc|cc}
A&B&A\to B&A\leftrightarrow B\\
\hline
F&F&T&T\\
F&T&T&F\\
T&F&F&F\\
T&T&T&T
\end{array}
\]
Two formulas \(\varphi,\psi\) are logically equivalent，写作 \(\varphi\equiv\psi\)，意思是它们在所有 truth assignments 下有相同 truth value。注意 \(\leftrightarrow\) 是 Boolean connective，而 \(\equiv\) 是 meta-level statement。

\begin{exerciseblock}[Lecture Notes Exercise 2]
用 truth tables 检查 Lecture Notes 中列出的 logical equivalences。
\end{exerciseblock}

\begin{solution}
典型等价式包括：
\[
A\land B\equiv B\land A,\qquad A\lor B\equiv B\lor A,
\]
\[
A\land(B\land C)\equiv(A\land B)\land C,\qquad
A\lor(B\lor C)\equiv(A\lor B)\lor C,
\]
\[
A\land(B\lor C)\equiv(A\land B)\lor(A\land C),
\]
\[
A\lor(B\land C)\equiv(A\lor B)\land(A\lor C),
\]
\[
\lnot(A\land B)\equiv\lnot A\lor\lnot B,\qquad
\lnot(A\lor B)\equiv\lnot A\land\lnot B,
\]
\[
A\to B\equiv\lnot A\lor B,\qquad
A\to B\equiv(\lnot B)\to(\lnot A).
\]
例如 de Morgan's law：
\[
\begin{array}{cc|c|c|c}
A&B&A\lor B&\lnot(A\lor B)&\lnot A\land\lnot B\\
\hline
F&F&F&T&T\\
F&T&T&F&F\\
T&F&T&F&F\\
T&T&T&F&F
\end{array}
\]
最后两列相同，所以 \(\lnot(A\lor B)\equiv\lnot A\land\lnot B\)。Contrapositive 也可直接检查：
\[
\begin{array}{cc|cc}
A&B&A\to B&(\lnot B)\to(\lnot A)\\
\hline
F&F&T&T\\
F&T&T&T\\
T&F&F&F\\
T&T&T&T
\end{array}
\]
\end{solution}

\begin{explanation}
证明 logical equivalence 的标准方法是比较 final columns。三变量 laws 需要八行 truth table；若两边每一行都相同，就说明它们定义同一个 truth function。
\end{explanation}

\section{Valid deductions}

一个 deduction valid 的意思是：不存在一个 truth assignment 使所有 premises 为 true，而 conclusion 为 false。

\begin{exerciseblock}[Lecture Notes Exercise 3]
检查
\[
A\to B,\quad B\to C,\quad \therefore A\to C
\]
以及
\[
\lnot A,\quad B\to A,\quad \therefore \lnot B.
\]
\end{exerciseblock}

\begin{solution}
第一条 valid。若 \(A=T\)，由 \(A\to B\) 得 \(B=T\)，再由 \(B\to C\) 得 \(C=T\)，所以 \(A\to C\) true。若 \(A=F\)，则 \(A\to C\) automatically true。

第二条 valid。若 conclusion \(\lnot B\) false，则 \(B=T\)。由 \(B\to A\) 得 \(A=T\)，但 premise \(\lnot A\) 要求 \(A=F\)，矛盾。
\end{solution}

\begin{explanation}
本题的解释是：先判断题目要求的是 translation, truth-table computation, validity test, or quantified statement，然后逐步展开对应 definition。这样的写法把每一个 conclusion 都连接到 truth value、counterexample row、witness, or universal instantiation，避免用直觉替代 formal reasoning。
\end{explanation}

\begin{exerciseblock}[Lecture Notes Exercise 4]
检查
\[
(A\lor B)\land\lnot B,\quad \therefore A.
\]
\end{exerciseblock}

\begin{solution}
\[
\begin{array}{cc|ccc}
A&B&A\lor B&\lnot B&(A\lor B)\land\lnot B\\
\hline
F&F&F&T&F\\
F&T&T&F&F\\
T&F&T&T&T\\
T&T&T&F&F
\end{array}
\]
premise true 的唯一 row 是 \(A=T,B=F\)，所以 conclusion \(A\) 必然 true。
\end{solution}

\begin{explanation}
本题的解释是：先判断题目要求的是 translation, truth-table computation, validity test, or quantified statement，然后逐步展开对应 definition。这样的写法把每一个 conclusion 都连接到 truth value、counterexample row、witness, or universal instantiation，避免用直觉替代 formal reasoning。
\end{explanation}

\begin{exerciseblock}[Lecture Notes Exercise 5]
证明
\[
A\to B,\qquad B\to C,\qquad C\to A
\]
等价于 \(A\leftrightarrow B\leftrightarrow C\)。
\end{exerciseblock}

\begin{solution}
为避免 chained biconditional 的 ambiguity，把题目理解为
\[
(A\leftrightarrow B)\land(B\leftrightarrow C),
\]
也就是 \(A,B,C\) have the same truth value。

若三条 implications 都 true，且 \(A=T\)，则 \(B=T\) 且 \(C=T\)。若 \(A=F\)，则 \(C\) 不能 true，否则 \(C\to A\) false；又 \(B\) 不能 true，否则 \(B\to C\) false。因此三者同时 true 或同时 false。反过来，若三者 truth value 相同，则三条 implications 都 true。因此两边 equivalent。
\end{solution}

\begin{explanation}
本题的解释是：先判断题目要求的是 translation, truth-table computation, validity test, or quantified statement，然后逐步展开对应 definition。这样的写法把每一个 conclusion 都连接到 truth value、counterexample row、witness, or universal instantiation，避免用直觉替代 formal reasoning。
\end{explanation}

\section{Tautology and contradiction}

Formula 若每一行都 true，称为 tautology；若每一行都 false，称为 contradiction；否则是 neither。例如 \(A\lor\lnot A\) 是 tautology，\(A\land\lnot A\) 是 contradiction。

\section{Worksheet 1: propositional logic in practice}

\begin{exerciseblock}[Worksheet 1 Exercise 1]
判断以下句子是否为 propositions。
\end{exerciseblock}

\begin{solution}
\begin{enumerate}[label=(\alph*)]
  \item \(2+2=4\) 是 proposition，true。
  \item \(x+1=3\) 不是 proposition，因为 \(x\) 没有指定。
  \item \(7\) is prime 是 proposition，true。
  \item ``Close the door'' 不是 proposition，因为它是 command。
  \item Every even number is divisible by \(2\) 是 proposition，true。
\end{enumerate}
\end{solution}

\begin{explanation}
本题的解释是：先判断题目要求的是 translation, truth-table computation, validity test, or quantified statement，然后逐步展开对应 definition。这样的写法把每一个 conclusion 都连接到 truth value、counterexample row、witness, or universal instantiation，避免用直觉替代 formal reasoning。
\end{explanation}

\begin{exerciseblock}[Worksheet 1 Exercise 2]
翻译 \(P\land Q\), \(P\lor Q\), \(\lnot P\), \(P\to Q\), \(P\leftrightarrow Q\)。
\end{exerciseblock}

\begin{solution}
\(P\land Q\)：\(P\) and \(Q\) are both true。\quad
\(P\lor Q\)：at least one of \(P,Q\) is true。\quad
\(\lnot P\)：\(P\) is false。\quad
\(P\to Q\)：if \(P\), then \(Q\)。\quad
\(P\leftrightarrow Q\)：\(P\) if and only if \(Q\)，即两者 truth value 相同。
\end{solution}

\begin{explanation}
本题的解释是：先判断题目要求的是 translation, truth-table computation, validity test, or quantified statement，然后逐步展开对应 definition。这样的写法把每一个 conclusion 都连接到 truth value、counterexample row、witness, or universal instantiation，避免用直觉替代 formal reasoning。
\end{explanation}

\begin{exerciseblock}[Worksheet 1 Exercise 3]
令 \(P=\) ``The homework is submitted''，\(Q=\) ``The homework is graded''。翻译四个 English statements。
\end{exerciseblock}

\begin{solution}
\[
P\land Q,\qquad P\land\lnot Q,\qquad Q\to P,\qquad Q\leftrightarrow P.
\]
第三个是 \(Q\to P\)，因为 ``If the homework is graded'' 是 antecedent。
\end{solution}

\begin{explanation}
本题的解释是：先判断题目要求的是 translation, truth-table computation, validity test, or quantified statement，然后逐步展开对应 definition。这样的写法把每一个 conclusion 都连接到 truth value、counterexample row、witness, or universal instantiation，避免用直觉替代 formal reasoning。
\end{explanation}

\begin{exerciseblock}[Worksheet 1 Exercise 4]
构造 \(\lnot P\lor Q\), \(P\to Q\), \((P\land Q)\to P\), \(P\to(Q\to P)\) 的 truth tables。
\end{exerciseblock}

\begin{solution}
\[
\begin{array}{cc|cccccc}
P&Q&\lnot P&\lnot P\lor Q&P\to Q&P\land Q&(P\land Q)\to P&P\to(Q\to P)\\
\hline
F&F&T&T&T&F&T&T\\
F&T&T&T&T&F&T&T\\
T&F&F&F&F&F&T&T\\
T&T&F&T&T&T&T&T
\end{array}
\]
\end{solution}

\begin{explanation}
本题的解释是：先判断题目要求的是 translation, truth-table computation, validity test, or quantified statement，然后逐步展开对应 definition。这样的写法把每一个 conclusion 都连接到 truth value、counterexample row、witness, or universal instantiation，避免用直觉替代 formal reasoning。
\end{explanation}

\begin{exerciseblock}[Worksheet 1 Exercise 5]
判断 \(P\lor\lnot P\), \(P\land\lnot P\), \((P\land Q)\to P\), \(P\to Q\)。
\end{exerciseblock}

\begin{solution}
\[
P\lor\lnot P\text{ is a tautology},\quad
P\land\lnot P\text{ is a contradiction},
\]
\[
(P\land Q)\to P\text{ is a tautology},\quad
P\to Q\text{ is neither}.
\]
理由是：\(P\lor\lnot P\) 每一行都有一个 side 为 true；\(P\land\lnot P\) 每一行都有一个 side 为 false。\((P\land Q)\to P\) 的 antecedent true 时 \(P\) 已经 true；\(P\to Q\) 在 \(P=T,Q=F\) 时 false，在其他 rows 中 true。
\end{solution}

\begin{explanation}
本题的解释是：先判断题目要求的是 translation, truth-table computation, validity test, or quantified statement，然后逐步展开对应 definition。这样的写法把每一个 conclusion 都连接到 truth value、counterexample row、witness, or universal instantiation，避免用直觉替代 formal reasoning。
\end{explanation}

\begin{exerciseblock}[Worksheet 1 Exercise 6]
令 \(P=\) ``It is raining''，\(Q=\) ``The ground is wet''。解释 \(P\to Q\)。
\end{exerciseblock}

\begin{solution}
\(P\to Q\) 表示 ``If it is raining, then the ground is wet''。它不 assert that it is raining，因为 implication 只说明 \(P\) true 时 \(Q\) 必须 true。若没有下雨，但地面因清洁而湿，则 \(P=F,Q=T\)，所以 \(P\to Q=T\)。
\end{solution}

\begin{explanation}
本题的解释是：先判断题目要求的是 translation, truth-table computation, validity test, or quantified statement，然后逐步展开对应 definition。这样的写法把每一个 conclusion 都连接到 truth value、counterexample row、witness, or universal instantiation，避免用直觉替代 formal reasoning。
\end{explanation}

\begin{exerciseblock}[Worksheet 1 Exercise 7]
说明 affirming the consequent 与 denying the antecedent invalid。
\end{exerciseblock}

\begin{solution}
\[
P\to Q,\quad Q,\quad \therefore P
\]
是 affirming the consequent。取 \(P=F,Q=T\)，premises true 但 conclusion false。
\[
P\to Q,\quad \lnot P,\quad \therefore \lnot Q
\]
是 denying the antecedent。同样取 \(P=F,Q=T\)，premises true 但 conclusion false。
\end{solution}

\begin{explanation}
本题的解释是：先判断题目要求的是 translation, truth-table computation, validity test, or quantified statement，然后逐步展开对应 definition。这样的写法把每一个 conclusion 都连接到 truth value、counterexample row、witness, or universal instantiation，避免用直觉替代 formal reasoning。
\end{explanation}

\begin{exerciseblock}[Worksheet 1 Exercise 8]
证明
\[
P\to Q\equiv\lnot P\lor Q,\quad
\lnot(P\lor Q)\equiv\lnot P\land\lnot Q,\quad
\lnot(P\land Q)\equiv\lnot P\lor\lnot Q.
\]
\end{exerciseblock}

\begin{solution}
第一条来自 implication truth table：两边只在 \(P=T,Q=F\) 时 false。后两条是 de Morgan's laws，可由四行 truth table 逐行检查。语义上，``not at least one'' 等于 ``both not''，而 ``not both'' 等于 ``at least one not''。
\end{solution}

\begin{explanation}
本题的解释是：先判断题目要求的是 translation, truth-table computation, validity test, or quantified statement，然后逐步展开对应 definition。这样的写法把每一个 conclusion 都连接到 truth value、counterexample row、witness, or universal instantiation，避免用直觉替代 formal reasoning。
\end{explanation}

\begin{exerciseblock}[Worksheet 1 Exercise 9]
Negate and simplify \(P\land Q\), \(P\lor Q\), \(P\to Q\), \(P\leftrightarrow Q\)。
\end{exerciseblock}

\begin{solution}
\[
\lnot(P\land Q)\equiv\lnot P\lor\lnot Q,\qquad
\lnot(P\lor Q)\equiv\lnot P\land\lnot Q,
\]
\[
\lnot(P\to Q)\equiv P\land\lnot Q,
\]
\[
\lnot(P\leftrightarrow Q)\equiv(P\land\lnot Q)\lor(\lnot P\land Q).
\]
这些 rewritten forms 都把 outer negation 消掉了。前两条使用 de Morgan's laws；第三条保留 implication false 的唯一 row；第四条表示 \(P,Q\) truth values different，也就是 exclusive-or pattern。
\end{solution}

\begin{explanation}
本题的解释是：先判断题目要求的是 translation, truth-table computation, validity test, or quantified statement，然后逐步展开对应 definition。这样的写法把每一个 conclusion 都连接到 truth value、counterexample row、witness, or universal instantiation，避免用直觉替代 formal reasoning。
\end{explanation}

\section{Homework 1: propositional logic}

\begin{exerciseblock}[Homework 1 Problem 1]
Insert parentheses and rewrite in words。
\end{exerciseblock}

\begin{solution}
按 standard precedence：
\[
(\lnot A\land B)\lor C,
\]
即 either \(A\) is false and \(B\) is true, or \(C\) is true。
\[
A\to(B\land\lnot C),
\]
即 if \(A\), then \(B\) and not \(C\)。
\[
\lnot A\leftrightarrow(B\lor C),
\]
即 \(A\) is false iff at least one of \(B,C\) is true。
\end{solution}

\begin{explanation}
本题的解释是：先判断题目要求的是 translation, truth-table computation, validity test, or quantified statement，然后逐步展开对应 definition。这样的写法把每一个 conclusion 都连接到 truth value、counterexample row、witness, or universal instantiation，避免用直觉替代 formal reasoning。
\end{explanation}

\begin{exerciseblock}[Homework 1 Problem 2]
Truth tables for \((A\lor B)\to C\), \(A\to(B\to A)\), \((A\land B)\leftrightarrow(B\land A)\)。
\end{exerciseblock}

\begin{solution}
\[
\begin{array}{ccc|ccccccc}
A&B&C&A\lor B&(A\lor B)\to C&B\to A&A\to(B\to A)&A\land B&B\land A&(A\land B)\leftrightarrow(B\land A)\\
\hline
F&F&F&F&T&T&T&F&F&T\\
F&F&T&F&T&T&T&F&F&T\\
F&T&F&T&F&F&T&F&F&T\\
F&T&T&T&T&F&T&F&F&T\\
T&F&F&T&F&T&T&F&F&T\\
T&F&T&T&T&T&T&F&F&T\\
T&T&F&T&F&T&T&T&T&T\\
T&T&T&T&T&T&T&T&T&T
\end{array}
\]
\end{solution}

\begin{explanation}
本题的解释是：先判断题目要求的是 translation, truth-table computation, validity test, or quantified statement，然后逐步展开对应 definition。这样的写法把每一个 conclusion 都连接到 truth value、counterexample row、witness, or universal instantiation，避免用直觉替代 formal reasoning。
\end{explanation}

\begin{exerciseblock}[Homework 1 Problem 3]
Classify \(A\lor\lnot A\), \(A\land\lnot A\), \((A\to B)\lor(B\to A)\)。
\end{exerciseblock}

\begin{solution}
\(A\lor\lnot A\) is a tautology；\(A\land\lnot A\) is a contradiction；\((A\to B)\lor(B\to A)\) is a tautology。最后一条中，两条 implications 不可能同时 false，因为第一条 false 要求 \(A=T,B=F\)，第二条 false 要求 \(B=T,A=F\)。
\end{solution}

\begin{explanation}
本题的解释是：先判断题目要求的是 translation, truth-table computation, validity test, or quantified statement，然后逐步展开对应 definition。这样的写法把每一个 conclusion 都连接到 truth value、counterexample row、witness, or universal instantiation，避免用直觉替代 formal reasoning。
\end{explanation}

\begin{exerciseblock}[Homework 1 Problem 4]
Show \(\lnot(A\lor B)\equiv\lnot A\land\lnot B\) and \(A\leftrightarrow B\equiv\lnot A\leftrightarrow\lnot B\)。
\end{exerciseblock}

\begin{solution}
第一条是 de Morgan's law，可由四行 truth table 得到。第二条中，\(A\leftrightarrow B\) 检查 \(A,B\) 是否同 truth value；同时 negating 两边不会改变 ``same'' 或 ``different''，所以
\[
A\leftrightarrow B\equiv\lnot A\leftrightarrow\lnot B.
\]
\end{solution}

\begin{explanation}
本题的解释是：先判断题目要求的是 translation, truth-table computation, validity test, or quantified statement，然后逐步展开对应 definition。这样的写法把每一个 conclusion 都连接到 truth value、counterexample row、witness, or universal instantiation，避免用直觉替代 formal reasoning。
\end{explanation}

\begin{exerciseblock}[Homework 1 Problem 5]
Rewrite \(\lnot(A\to B)\) without \(\to\)。
\end{exerciseblock}

\begin{solution}
\[
\lnot(A\to B)\equiv\lnot(\lnot A\lor B)\equiv A\land\lnot B.
\]
换句话说，否定 implication 时要描述它失败的唯一情况：antecedent \(A\) true，但 consequent \(B\) false。因此 final formula 不需要 \(\to\)，只需要 conjunction 和 negation。
\end{solution}

\begin{explanation}
本题的解释是：先判断题目要求的是 translation, truth-table computation, validity test, or quantified statement，然后逐步展开对应 definition。这样的写法把每一个 conclusion 都连接到 truth value、counterexample row、witness, or universal instantiation，避免用直觉替代 formal reasoning。
\end{explanation}

\begin{exerciseblock}[Homework 1 Problem 6]
Decide valid deductions。
\end{exerciseblock}

\begin{solution}
\begin{enumerate}[label=(\alph*)]
  \item \(A\to B,\ B\to C,\ \therefore A\to C\) is valid。
  \item \(A\lor B,\ \lnot A,\ \therefore B\) is valid。
  \item \(A\to B,\ \lnot A,\ \therefore \lnot B\) is invalid。Counterexample：\(A=F,B=T\)。
\end{enumerate}
\end{solution}

\begin{explanation}
本题的解释是：先判断题目要求的是 translation, truth-table computation, validity test, or quantified statement，然后逐步展开对应 definition。这样的写法把每一个 conclusion 都连接到 truth value、counterexample row、witness, or universal instantiation，避免用直觉替代 formal reasoning。
\end{explanation}

\section{Predicate logic}

Predicate logic 允许 variables、predicates 和 quantifiers。Domain of discourse 必须明确。若 \(H(x)\) 表示 ``\(x\) is human''，\(M(x)\) 表示 ``\(x\) is mortal''，则
\[
\forall x(H(x)\to M(x)),\qquad H(s),\qquad \therefore M(s)
\]
表达 ``All humans are mortal; Confucius is human; therefore Confucius is mortal''。

\begin{exerciseblock}[Lecture Notes Exercise 6]
Translate three English sentences into predicate formulas。
\end{exerciseblock}

\begin{solution}
令 \(C(x)\) 表示 \(x\) is a couple，\(F(x)\) 表示 \(x\) has fights：
\[
\forall x(C(x)\to F(x)).
\]
令 \(B(y)\) 表示 \(y\) is a book，\(R(x,y)\) 表示 \(x\) is reading \(y\)：
\[
\forall x\exists y\exists z\bigl(B(y)\land B(z)\land y\ne z\land R(x,y)\land R(x,z)\bigr).
\]
令 \(c\) 表示 my classmate，\(Fr(c,x)\) 表示 \(x\) is a friend of \(c\)：
\[
\exists x\bigl(Fr(c,x)\land\forall y(Fr(c,y)\to y=x)\bigr),
\]
也可写成 \(\exists!x\,Fr(c,x)\)。
\end{solution}

\begin{explanation}
本题的解释是：先判断题目要求的是 translation, truth-table computation, validity test, or quantified statement，然后逐步展开对应 definition。这样的写法把每一个 conclusion 都连接到 truth value、counterexample row、witness, or universal instantiation，避免用直觉替代 formal reasoning。
\end{explanation}

\begin{exerciseblock}[Lecture Notes Exercise 7]
Rewrite \(\lnot\forall xP(x)\) using \(\exists\)。
\end{exerciseblock}

\begin{solution}
\[
\lnot\forall xP(x)\equiv\exists x\lnot P(x).
\]
意思是：not every \(x\) has property \(P\)，等价于 there exists at least one \(x\) that does not have property \(P\)。
\end{solution}

\begin{explanation}
本题的解释是：先判断题目要求的是 translation, truth-table computation, validity test, or quantified statement，然后逐步展开对应 definition。这样的写法把每一个 conclusion 都连接到 truth value、counterexample row、witness, or universal instantiation，避免用直觉替代 formal reasoning。
\end{explanation}

\begin{exerciseblock}[Lecture Notes Exercise 8]
Give examples of laws of predicate logic。
\end{exerciseblock}

\begin{solution}
常用 laws 包括：
\[
\forall x(P(x)\land Q(x))\equiv(\forall xP(x))\land(\forall xQ(x)),
\]
\[
\exists x(P(x)\lor Q(x))\equiv(\exists xP(x))\lor(\exists xQ(x)),
\]
\[
\lnot\forall xP(x)\equiv\exists x\lnot P(x),\qquad
\lnot\exists xP(x)\equiv\forall x\lnot P(x),
\]
\[
\forall x\forall yP(x,y)\equiv\forall y\forall xP(x,y),\qquad
\exists x\exists yP(x,y)\equiv\exists y\exists xP(x,y).
\]
\end{solution}

\begin{explanation}
本题的解释是：先判断题目要求的是 translation, truth-table computation, validity test, or quantified statement，然后逐步展开对应 definition。这样的写法把每一个 conclusion 都连接到 truth value、counterexample row、witness, or universal instantiation，避免用直觉替代 formal reasoning。
\end{explanation}

\begin{exerciseblock}[Lecture Notes Exercise 9]
Translate and compare
\[
\forall x\exists y(x<y),\qquad \exists y\forall x(x<y).
\]
\end{exerciseblock}

\begin{solution}
第一句：for every \(x\)，there exists a \(y\) with \(x<y\)。在 \(\N,\Z,\Q,\R\) 等 usual ordered sets 中，这表示每个 element 上方还有更大的 element。

第二句：there exists one \(y\) such that every \(x\) satisfies \(x<y\)。这表示存在一个 single element 严格大于所有 \(x\)，在 usual \(\N,\Z,\Q,\R\) 中 false。
\end{solution}

\begin{explanation}
Mixed quantifiers 的 order 不能随便交换。\(\forall x\exists y\) 允许 \(y\) depend on \(x\)；\(\exists y\forall x\) 要求同一个 \(y\) 同时服务所有 \(x\)。
\end{explanation}

\begin{exerciseblock}[Lecture Notes Exercise 10]
List variants of the listed equivalences which fail。
\end{exerciseblock}

\begin{solution}
以下 variants 一般 fail：
\[
\forall x(P(x)\lor Q(x))\not\equiv(\forall xP(x))\lor(\forall xQ(x)).
\]
Countermodel：domain \(\{1,2\}\)，令 \(P(1)\) true, \(P(2)\) false, \(Q(1)\) false, \(Q(2)\) true。则每个 \(x\) 都满足 \(P(x)\lor Q(x)\)，但 neither \(\forall xP(x)\) nor \(\forall xQ(x)\)。
\[
\exists x(P(x)\land Q(x))\not\equiv(\exists xP(x))\land(\exists xQ(x)).
\]
同一 countermodel 下右边 true，但没有同一个 \(x\) 同时满足 \(P,Q\)。
\[
\forall x\exists yP(x,y)\not\equiv\exists y\forall xP(x,y).
\]
取 domain \(\N\)，令 \(P(x,y)\) 为 \(x<y\)，左边 true，右边 false。
\end{solution}

\begin{explanation}
本题的解释是：先判断题目要求的是 translation, truth-table computation, validity test, or quantified statement，然后逐步展开对应 definition。这样的写法把每一个 conclusion 都连接到 truth value、counterexample row、witness, or universal instantiation，避免用直觉替代 formal reasoning。
\end{explanation}

\section{Homework 2: predicate logic and quantifiers}

\begin{exerciseblock}[Homework 2 Problem 1]
Domain 是 all students at a university，\(S(x)\): \(x\) studies mathematics，\(P(x)\): \(x\) passed the exam。Translate。
\end{exerciseblock}

\begin{solution}
\[
\forall x(S(x)\to P(x)),\qquad
\exists xP(x),\qquad
\exists x(S(x)\land\lnot P(x)).
\]
第一句说每个 mathematics student 都 passed the exam；第二句只要求至少一个 student passed；第三句要求同一个 student 同时 studies mathematics 且 did not pass。Domain 已经是 all students，所以不需要另写 student predicate。
\end{solution}

\begin{explanation}
本题的解释是：先判断题目要求的是 translation, truth-table computation, validity test, or quantified statement，然后逐步展开对应 definition。这样的写法把每一个 conclusion 都连接到 truth value、counterexample row、witness, or universal instantiation，避免用直觉替代 formal reasoning。
\end{explanation}

\begin{exerciseblock}[Homework 2 Problem 2]
Translate into English：
\[
\forall x\exists y(x\ne y),\qquad \exists y\forall x(x\le y).
\]
\end{exerciseblock}

\begin{solution}
第一句：for every object \(x\)，there is an object \(y\) different from \(x\)。这等价于 domain 至少有两个 elements。

第二句：there exists an object \(y\) such that every \(x\) is less than or equal to \(y\)。在 ordered domain 中，这表示 \(y\) is a greatest element。
\end{solution}

\begin{explanation}
本题的解释是：先判断题目要求的是 translation, truth-table computation, validity test, or quantified statement，然后逐步展开对应 definition。这样的写法把每一个 conclusion 都连接到 truth value、counterexample row、witness, or universal instantiation，避免用直觉替代 formal reasoning。
\end{explanation}

\begin{exerciseblock}[Homework 2 Problem 3]
Negate and move negation inside quantifiers。
\end{exerciseblock}

\begin{solution}
\[
\lnot\forall xP(x)\equiv\exists x\lnot P(x),
\]
\[
\lnot\exists x(P(x)\land Q(x))\equiv\forall x\lnot(P(x)\land Q(x))
\equiv\forall x(\lnot P(x)\lor\lnot Q(x)).
\]
第一条把 ``not all'' 改写为 ``there exists a counterexample''。第二条先把 \(\lnot\exists\) 改成 \(\forall\lnot\)，再用 de Morgan's law 把 conjunction 的 negation 推到 predicates 上。
\end{solution}

\begin{explanation}
本题的解释是：先判断题目要求的是 translation, truth-table computation, validity test, or quantified statement，然后逐步展开对应 definition。这样的写法把每一个 conclusion 都连接到 truth value、counterexample row、witness, or universal instantiation，避免用直觉替代 formal reasoning。
\end{explanation}

\begin{exerciseblock}[Homework 2 Problem 4]
Determine whether the argument is logically valid：
\[
\forall x(P(x)\to Q(x)),\qquad \exists xP(x),\qquad \therefore \exists xQ(x).
\]
\end{exerciseblock}

\begin{solution}
Valid。由 \(\exists xP(x)\)，取 witness \(a\) 使 \(P(a)\) true。由 \(\forall x(P(x)\to Q(x))\)，对 \(a\) 做 universal instantiation，得 \(P(a)\to Q(a)\)。于是 \(Q(a)\) true，因此 \(\exists xQ(x)\) true。
\end{solution}

\begin{explanation}
这题的关键是 witness reasoning：existential premise 给出某个 \(a\)，universal premise 可以用于同一个 \(a\)。不能把 witness 随意换成另一个 object。
\end{explanation}
